\chapter{Introdução}
\label{ch:introducao}

\section{Conceitos Gerais}
\label{sec:conceitos}

Atualmente, a maioria dos usuários interage com uma grande variedade de sistemas computacionais, tanto no trabalho quanto em casa. Entre os sistemas utilizados, podem-se citar máquinas fotográficas digitais, telefones celulares, caixas eletrônicos, aparelhos de DVD, reprodutores de música portáteis e {\em tablets}. No entanto, muitos dos usuários não percebem a variedade e a quantidade de sistemas utilizados diariamente.

Cada vez mais, existe a tendência do aumento da interconectividade entre estes sistemas, devido ao fato da internet estar cada vez mais presente e disponível em qualquer lugar. Assim, pesquisas recentes preveem a integração completa de todos estes sistemas em uma federação de sistemas computacionais, que podem ser descritos como Ambientes Inteligentes \cite{heckmann05}.

Uma visão para o conceito de ambientes inteligentes foi discutida por \citetexto{weiser91} em seu artigo seminal, onde ele apresentou um cenário em que o computador se tornaria invisível e indistinguível de seu ambiente. Neste cenário, a interação entre humanos e computadores ocorreria não como a usual -- onde o humano dedica 100\% de sua atenção ao computador --, mas constituíra-se de pequenas interações rápidas, geralmente em locais e dispositivos distintos, muitas vezes utilizando a periferia da atenção do usuário \cite{weiser96} ou sem ele se dar conta de estar interagindo com um sistema computacional.

Para que este cenário seja possível, é necessário que os aplicativos tenham condições de considerar o contexto no qual se encontram e adaptar-se a este contexto, ou seja, é necessário que sejam ``sensíveis a contexto''. Embora existam várias definições de contexto, uma das mais utilizadas é a de que um contexto é qualquer informação que pode ser usada para caracterizar a situação de uma entidade, onde uma entidade pode ser uma pessoa, local ou objeto considerado relevante à interação entre o usuário e a aplicação \cite{dey2001}.

Para que um aplicativo possa se adaptar a um usuário de forma eficaz, é importante que ele conheça não apenas o histórico de contextos visitados por um usuário, mas conheça quem o usuário é. Atualmente, muitos trabalhos têm buscado maneiras de coletar informações a respeito dos usuários, com o objetivo de conhecer suas preferências, possibilitando a oferta de serviços personalizados. Assim, estas aplicações organizam seu conhecimento a respeito do usuário em perfis de usuário, os quais armazenam as informações em um formato baseado em um modelo de usuário específico \cite{viviani10}.

Pesquisas nesta área têm indicado que não apenas é interessante conhecer o contexto do usuário no momento em que um determinado serviço está sendo oferecido, mas que o histórico de contextos visitados pode ser uma informação valiosa na adaptação da aplicação ao usuário. Este histórico de contextos recebe o nome de Trilha \cite{silva10}.


%%%%%%%%%%%%%%%%%%%%%%%%%%%%%%%%%%%%%%%%%%%%%%%%%%%%%%%%%%%%%%%%%%%%

\section{Motivação}
\label{sec:motivacao}

Desde que os primeiros computadores se tornaram disponíveis, tem havido uma tendência: o custo e o tamanho dos dispositivos têm diminuído, enquanto sua presença têm se tornado cada vem mais ubíqua. E com o advento da internet, cada vez mais serviços tem sido disponibilizados em qualquer lugar que o usuário esteja. Estes aspectos deram origem a uma visão onde serviços computacionais estão disponíveis a qualquer momento e em qualquer lugar, denominada de Computação Ubíqua \cite{weiser91}.

\citetexto{weiser96} discutem as consequências futuras, caso esta tendência continue -- e não há nada que demonstre que ela irá reduzir. Caso o modelo atual de interação humano-computador seja mantido, é provável que cada vez mais os usuários estejam sobrecarregados de interações computacionais e informações, causando assim detrimento na qualidade de vida destes. Por isso, os autores apresentam o conceito de Tecnologia Calma, onde a interação entre os computadores e o humano ocorrem na periferia de atenção do usuário. Neste conceito, o computador desaparece da consciência do usuário, a não ser no momento em que seja necessário, quando então chama atenção para si.

\citetexto{satyanarayanan01} chama a atenção para o fato de que um sistema só pode se tornar invisível se o mesmo for pró-ativo, e só pode se tornar pró-ativo se souber a intenção do usuário. Muitas pesquisas têm sido feitas sobre a melhor forma de se obter a intenção do usuário e compreender quem o usuário é, o que ele está tentando realizar, e como o sistema pode recomendar opções a ele ou adaptar-se a ele.

Para que um sistema tenha sucesso em obter a intenção do usuário, é necessário que ele conheça o usuário, o contexto que o usuário está inserido e as ações que ele tomou anteriormente quando se encontrava no mesmo contexto ou em contextos semelhantes. As informações do usuário podem se enquadrar em um Perfil de Usuário, enquanto as informações históricas de contexto se encaixam na ideia de Trilha.

Embora existam alguns trabalhos que busquem inferir conhecimento a respeito de um usuário a partir de sua trilha, os trabalhos anteriores operam dentro de um contexto específico (como Educação \cite{barbosa11} ou Comércio \cite{franco11}) ou trabalham com o modelo de usuário de forma estática.

Neste trabalho, o conceito de Entidade é usado no lugar do conceito clássico de Usuário. Enquanto geralmente o usuário compreende a pessoa que está utilizando o aplicativo, uma entidade refere-se a qualquer objeto (físico ou não) capaz de controlar um aplicativo. Desta forma, uma entidade pode ser uma pessoa, um objeto móvel (como um carro) ou imóvel (como uma sala), ou até mesmo um sistema.

%%%%%%%%%%%%%%%%%%%%%%%%%%%%%%%%%%%%%%%%%%%%%%%%%%%%%%%%%%%%%%%%%%%%

\section{Objetivos}
\label{sec:objetivos}

O objetivo geral do trabalho é especificar, implementar e validar um modelo de gerenciamento de perfis chamado {\bf eProfile}, através do qual a trilha de uma entidade possa ser processada e transformada em um perfil de entidade. O modelo a ser apresentado é um modelo genérico, ou seja, um modelo que permite trabalhar com qualquer formato de trilha e gerar qualquer perfil de entidade requerido pelas aplicações.

Assim sendo, são objetivos específicos deste trabalho:

\begin{itemize}

  \item apresentar o estado-da-arte em gerenciamento de perfis;
  \item especificar o eProfile;
  \item desenvolver um protótipo para o modelo;
  \item validar o modelo a partir do protótipo apresentado.

\end{itemize}

%%%%%%%%%%%%%%%%%%%%%%%%%%%%%%%%%%%%%%%%%%%%%%%%%%%%%%%%%%%%%%%%%%%%

\section{Metodologia}
\label{sec:metodologia}

A viabilização dos objetivos propostos por este trabalho se dá através de uma metodologia cujo primeiro passo é o estudo de modelos já propostos para suporte ao gerenciamento de perfis. O objetivo deste primeiro passo é verificar como os modelos atuais gerenciam perfis, e quais são os desafios e limitações da área.

O segundo passo será a especificação do eProfile. O objetivo é determinar quais módulos são necessários para o modelo e como eles irão interagir entre si e com serviços externos.

A elaboração de um cenário para aplicação do modelo será o passo seguinte. Esta atividade tem por objetivo explorar os possíveis ganhos que serão trazidos pelo modelo, para que o mesmo possa ser ajustado conforme necessário.

Para validar o modelo proposto será desenvolvido um protótipo. Este desenvolvimento irá cobrir todos os componentes que se fizerem necessários para viabilizar o cenário descritos no passo anterior.

Por fim, será feita uma avaliação do modelo, onde será realizada uma simulação da operação de entidades, gerando e analisando trilhas de entidades e inferindo perfis através destas trilhas. Os dados desse experimento serão medidos e tabulados a fim de se verificar o desempenho do modelo.

%%%%%%%%%%%%%%%%%%%%%%%%%%%%%%%%%%%%%%%%%%%%%%%%%%%%%%%%%%%%%%%%%%%%

\section{Organização da Dissertação}
\label{sec:organizacao}

Este trabalho está organizado da seguinte forma. O Capítulo~\ref{ch:conceitos} descreve os principais conceitos referentes às áreas relacionadas a este trabalho. O Capítulo~\ref{ch:relacionados} descreve trabalhos relacionados na área de gerenciamento de perfis, apresentando os principais desafios da área. No Capítulo~\ref{ch:modelo} é proposto o modelo do eProfile, bem como sua forma de operação. O Capítulo~\ref{ch:implementacao} apresenta os aspectos de implementação e avaliação do modelo, e o Capítulo~\ref{ch:avaliacao} apresenta o cenário que serve de avaliação do modelo. Por fim, o Capítulo~\ref{ch:consideracoes} apresenta as considerações finais.
